\documentclass[a4paper, 12pt]{article}  %set doc type and font size
\usepackage{fontspec}

% math tools
\usepackage{amsmath}
\usepackage{amssymb}
\usepackage{amsfonts}

% set chinese
\usepackage{xeCJK}
\xeCJKsetup{AutoFakeBold=true, AutoFakeSlant=true}
\setmainfont{Times New Roman}
\setCJKmainfont{標楷體}
\XeTeXlinebreaklocale "zh"
\XeTeXlinebreakskip = 0pt plus 1pt

\title{Chapter 2　Homework Assignment}
\author{黃丞暘　611211106}
\date{}

\begin{document}
\maketitle % 顯示標題區塊

\section*{Problem 1: Transformations and Standardization}
Let $X$ be a continuous random variable.
    \begin{enumerate}
    \item Prove $Var(aX+b)=a^{2}Var(X)$.

    Sol.
    \begin{align*}
    Var(aX+b)
    &=E[(aX+b)^2]-E[(aX+b)]^2 \\
    &=E[a^2X^2+2abX+b^2]-[aE(X)+b]^2 \\
    &=a^2E[X^2]+2abE[X]+b^2-(a^2E[X]^2+2abE[X]+b^2) \\
    &=a^2(E[X^2]-E[X]^2) \\
    &=a^2Var(X)
    \end{align*}

    \item Let $Y=\frac{X-\mu}{\sigma}$. Show that ${E}[Y]=0$ and $Var(Y)=1$.

    Sol.

    Let $\mu=E[X],\ \sigma=\sqrt{Var(X)}>0$, therefore
    $$E[Y]=E\!\left[\frac{X-\mu}{\sigma}\right]=\frac{1}{\sigma}\big(E[X]-\mu\big)=\frac{1}{\sigma}(\mu-\mu)=0,$$
    $$Var(Y)=Var(\frac{1}{\sigma}X-\frac{\mu}{\sigma})=\left(\frac{1}{\sigma}\right)^2Var(X-\mu)=\frac{1}{\sigma^2}Var(X)=\frac{1}{\sigma^2}\sigma^2=1.$$

    \item If $X\sim \mathrm{Exp}(\lambda)$  derive the CDF and verify the memoryless property.

    Sol.

    CDF: $$F_X(x)=P(X\le x)=\int_{0}^{x}\lambda e^{-\lambda t}\,dt =\left[-e^{-\lambda t}\right]_{0}^{x}=1-e^{-\lambda x}.$$
    Thus
    $$
    F_X(x)=
    \begin{cases}
        0, & x<0 \\
        1-e^{-\lambda x}, & x\ge 0
    \end{cases}.
    $$
    Let $s,t\ge 0$. Then
    $$P(X>s+t\mid X>s)=\frac{P(X>s+t)}{P(X>s)}=\frac{e^{-\lambda(s+t)}}{e^{-\lambda s}}=e^{-\lambda t}=P(X>t).$$
    Hence the exponential distribution is memoryless.

    \end{enumerate}

    \textbf{SPC Link}: Why is standardization fundamental in control limits?

    Sol.

    在 SPC 中，標準化重要的原因是：
    控制界線是建立在機率的基礎上，
    透過 $\frac{X-\mu}{\sigma}$ 的轉換，我們可以把統計量的中心移至 0 ，並以標準差為距離作為統一的單位，
    使不同製程資料都能在相同的標準下判斷是否異常，
    這也是為甚麼常見的 3-$\sigma$ 控制界線具有意義。

\section*{Problem 2: Normal Sampling Theory}
Let $X_1,\dots,X_n \stackrel{iid}{\sim} N(\mu,\sigma^2)$
    \begin{enumerate}
        \item Derive the distribution of $\bar{X}$.
        
        Sol.

        $\bar{X}=\frac{\sum_{i=1}^n X_i}{n}=\frac{X_1+X_2+\dots +X_n}{n}$, then 
        $$E[\bar{X}]=\frac{E[X_1+X_2+\dots +X_n]}{n}=\frac{n\mu}{n}=\mu,$$
        $$Var(\bar{X})=\frac{Var[X_1+X_2+\dots +X_n]}{n^2}=\frac{n\sigma^2}{n^2}=\frac{\sigma^2}{n}.$$
        Hence $$\bar{X}\sim N\left(\mu,\frac{\sigma^2}{n}\right).$$

        \item Show $\dfrac{(n-1)S^2}{\sigma^2} \sim \chi^2_{n-1}$.
        
        Sol.
        
        Let $S^2=\frac{\sum_{i=1}^n(X_i-\bar X)^2}{n-1}$, and we already know the following:

        (1) $\bar X$ and $S^2$ are independent by Basu's theorem.

        (2) If $Z=\frac{X-\mu}{\sigma}\sim N(0,1)$, then $Z^2\sim \chi^2_1$.
        
        (3) If $Y_i\stackrel{iid}{\sim}\chi^2_1$, then $\sum_{i=1}^n Y_i\sim \chi^2_n$,　$i=1,\dots,n$.
        
        (4) If $Q\sim \chi^2_k$, then $M_Q(t)=(1-2t)^{-k/2}$ for $t<\frac{1}{2}$.
        
        Compute
        \begin{align*}
        \sum_{i=1}^n (X_i-\mu)^2
        &=\sum_{i=1}^n((X_i-\bar X)+(\bar X-\mu))^2 \\
        &=\sum_{i=1}^n (X_i-\bar X)^2 + n(\bar X-\mu)^2.
        \end{align*}

        Then divide by $\sigma^2$:
        $$\sum_{i=1}^n\left(\frac{X_i-\mu}{\sigma}\right)^2=\sum_{i=1}^n\left(\frac{X_i-\bar X}{\sigma}\right)^2+\left(\frac{\bar X-\mu}{\sigma /\sqrt n}\right)^2.$$
        Let
        $$U=\sum_{i=1}^n\left(\frac{X_i-\mu}{\sigma}\right)^2,　V=\sum_{i=1}^n\left(\frac{X_i-\bar X}{\sigma}\right)^2,　W=\left(\frac{\bar X-\mu}{\sigma /\sqrt n}\right)^2.$$
        Then $U=V+W$, and we can know
        $$U=\sum_{i=1}^n Z_i^2 \sim \chi^2_n,　V=\frac{1}{\sigma^2}\sum_{i=1}^n (X_i-\bar X)^2=\frac{(n-1)S^2}{\sigma^2},$$
        also $\bar X\sim N(\mu,\sigma^2/n)$, hence $$\frac{\bar X-\mu}{\sigma /\sqrt n}\sim N(0,1)\Rightarrow W\sim\chi^2_1.$$
        By (1) and $W$ is a function of $\bar X$ while
        $V=\frac{(n-1)S^2}{\sigma^2}$ is a function of $S^2$, we have $$V\perp \!\!\! \perp W.$$
        From $U=V+W$ and independence, we have $$M_U(t)=M_V(t)\,M_W(t).$$
        By (4), we have $$(1-2t)^{-n/2}=M_U(t)=M_V(t)\,(1-2t)^{-1/2}\Rightarrow M_V(t)=(1-2t)^{-(n-1)/2}.$$
        Since $M_V(t)$ is the mgf of $\chi^2_{n-1}$, finally, $$V=\frac{(n-1)S^2}{\sigma^2}\sim \chi^2_{n-1}.$$

        \item Deduce $T=\frac{\bar{X}-\mu}{S/\sqrt{n}} \sim t_{n-1}$.
        
        Sol.
        
        Let $Z=\frac{\bar X-\mu}{\sigma /\sqrt n}\sim N(0,1)$. From Q2, we know and let $U=\frac{(n-1)S^2}{\sigma^2}\sim \chi^2_{n-1}.$

        Compute
        \begin{align*}
        T=\frac{\bar X-\mu}{S/\sqrt n}=\frac{(\bar X-\mu)/(\sigma /\sqrt n)}{(S/\sqrt n)/(\sigma /\sqrt n)}=\frac{Z}{S/\sigma}
        &=\frac{Z}{\sqrt{(S/\sigma)^2}} \\
        &=\frac{Z}{\sqrt{(n-1)S^2/(n-1)\sigma^2}} \\
        &=\frac{Z}{\sqrt{U/(n-1)}}
        \end{align*}
        Because of the definition of $t$-distribution：if $Z\sim N(0,1)$、$U\sim\chi^2_\nu$ and $Z\perp \!\!\! \perp U$，then $$\frac{Z}{\sqrt{U/\nu}}\sim t_\nu.$$
        In conclusion, $$T=\frac{\bar X-\mu}{S/\sqrt n}=\frac{Z}{\sqrt{U/(n-1)}}\sim t_{n-1}$$

    \end{enumerate}

    \textbf{SPC Link}: Known vs unknown σ changes limits.

    Sol.

    在 SPC 中，若製程資料的標準差 $\sigma$ 已知，則 $\bar{X}$ 的控制界線會服從常態分配 $Z$ ；若 $\sigma$ 未知，則會以樣本標準差 $S$ 取代 $\sigma$ 來估計，因此統計量會服從 $t$ 分配，
    因此控制界線需要調整。

\section*{Problem 3: Discrete Structure and Limits}
    \begin{enumerate}
        \item Derive $E[X]$, $Var(X)$ for $X \sim \mathrm{Bin}(n,p)$.
        
        Sol.

        Let $K_i\stackrel{iid}{\sim}\mathrm{Bernoulli}(p)$, we compute
        $$E[K_i]=1\times p+0\times (1-p)=p,$$
        $$E[K_i]^2=1^2\times p+0^2\times (1-p)=p,$$
        $$Var(K_i)=p-p^2=p(1-p).$$
        
        Since $X=\sum_{i=1}^n K_i$. Therefore,
        $$E[X]=E\!\left[\sum_{i=1}^n K_i\right]=\sum_{i=1}^n E[K_i]=\sum_{i=1}^n p=np,$$
        $$Var(X)=Var\left(\sum_{i=1}^n K_i\right)=\sum_{i=1}^n Var(K_i)=\sum_{i=1}^n p(1-p)=np(1-p).$$
        
        \item Prove the Poisson limit of Binomial.
        
        Sol.

        Let $X_n\sim \mathrm{Bin}(n,p)$ with $np\to \lambda >0$ as $n\to\infty$ and $p\to 0$.
        $$P(X=k)=\binom{n}{k}p^k(1-p)^{n-k}.$$
        We compute
        $$\binom{n}{k}p^k=\frac{n(n-1)\cdots(n-k+1)}{k!}p^k=\frac{1}{k!}\left(\prod_{j=0}^{k-1}\frac{n-j}{n}\right)(np)^k\to\frac{1}{k!}\lambda^k,　\text{as}　　n\to\infty,$$
        $$(1-p)^{n-k}=(1-p)^n(1-p)^{-k}=\left[\left(1-p\right)^{1/p}\right]^{np}(1-p)^{-k}\to e^{-\lambda}\cdot 1,　\text{as}　　n\to\infty.$$
        In conclusion, $$P(X=k)=\binom{n}{k}p^k(1-p)^{n-k}\to \frac{e^{-\lambda}\lambda^k}{k!},　\text{as}　　n\to\infty.$$

        \item Give an SPC context favoring Poisson modeling.
        
        Sol.
        
        在統計製程管制當中，Poisson 分配最常被用來描述在固定時間或固定區域內，某事件發生的次數

        例如，在半導體製程產線中，工程師檢查每片晶片，並記錄晶片表面上的缺陷數量。

        令 $X$ 為每片晶片的缺陷數量，若缺陷在晶片表上隨機且彼此間獨立，且在固定區域內的缺陷發生數近似於常數，
        令 $\lambda$ 為每片晶片的平均缺陷數，
        則 $$X \sim \mathrm{Poisson}(\lambda)\Rightarrow P(X=x)=\frac{e^{-\lambda}\lambda^x}{x!},\ x=0,1,2,\dots$$
        
    \end{enumerate}

\section*{Problem 4: Mixtures and Overdispersion}
    \begin{enumerate}
        \item Show Poisson--Gamma $\Rightarrow$ Negative Binomial.
        
        Sol.

        Let $X\mid \Lambda=\lambda \sim \mathrm{Pois}(\lambda)$ and $\Lambda \sim \mathrm{Gamma}(\alpha,\beta)$ for $\alpha,　\beta>0$, i.e.
        $$f_\Lambda(\lambda)=\frac{e^{-\beta\lambda}\beta^\alpha}{\Gamma(\alpha)}\lambda^{\alpha-1},　\lambda>0.$$
        Then $$f_{X\mid \Lambda}(x\mid \lambda)=\frac{e^{-\lambda}\lambda^x}{x!},　 x=0,1,2,\dots$$
        So we can compute
        \begin{align*}
        P(X=x)
        &=\int_0^\infty f_{X\mid \Lambda}(x\mid \lambda)f_\Lambda(\lambda)\,d\lambda \\
        &=\int_0^\infty \frac{e^{-\lambda}\lambda^x}{x!}\cdot
        \frac{e^{-\beta\lambda}\beta^\alpha}{\Gamma(\alpha)}\lambda^{\alpha-1}d\lambda \\
        &=\frac{\beta^\alpha}{x!\,\Gamma(\alpha)}
        \int_0^\infty \lambda^{x+\alpha-1}e^{-(\beta+1)\lambda}\,d\lambda.
        \end{align*}
        Use $\int_0^\infty \lambda^{t-1}e^{-k\lambda}\,d\lambda=\frac{\Gamma(t)}{k^t}$ with $t=x+\alpha$ and $k=\beta+1$:
        \begin{align*}
            P(X=x)
            &=\frac{\beta^\alpha}{x!\Gamma(\alpha)}\cdot \frac{\Gamma(x+\alpha)}{(\beta+1)^{x+\alpha}} \\
            &=\frac{\Gamma(x+\alpha)}{\Gamma(\alpha)x!}\cdot \left(\frac{\beta}{\beta+1}\right)^\alpha \left(\frac{1}{\beta+1}\right)^x.
        \end{align*}
        Let
        $$p=\frac{\beta}{\beta+1},　1-p=\frac{1}{\beta+1},　p\in[0, 1].$$
        Therefore $$P(X=x)=\frac{\Gamma(x+\alpha)}{\Gamma(\alpha)x!}p^{\alpha}(1-p)^x,　x=0,1,2,\dots$$
        It is the pmf of a Negative Binomial distribution: $$X\sim \mathrm{NegBin}(\alpha,p)　\text{for number of failures}　x　\text{before}　\alpha　\text{successes}$$

        \item Show Beta--Binomial via marginalization.
        
        Sol.

        Let $$X\mid P=p \sim \mathrm{Bin}(n,p),　P\sim \mathrm{Beta}(\alpha,\beta),　\alpha,\beta>0,$$
        where$$f_P(p)=\frac{1}{B(\alpha, \beta)}p^{\alpha-1}(1-p)^{\beta-1},　0<p<1,$$
        and $$B(\alpha,\beta)=\frac{\Gamma(\alpha)\Gamma(\beta)}{\Gamma(\alpha+\beta)}.$$
        Also,$$P_{X\mid P}(x\mid p)=\binom{n}{x}p^x(1-p)^{n-x},　x=0,1,\dots,n.$$
        Then we compute
        \begin{align*}
        P(X=x)
        &=\int_0^1 P_{X\mid P}(x\mid p)f_P(p)dp \\
        &=\int_0^1 \binom{n}{x}p^x(1-p)^{n-x}\cdot \frac{1}{B(\alpha,\beta)}p^{\alpha-1}(1-p)^{\beta-1}dp \\
        &=\binom{n}{x}\frac{1}{B(\alpha,\beta)}\int_0^1 p^{x+\alpha-1}(1-p)^{n-x+\beta-1}dp \\
        &=\binom{n}{x}\frac{B(x+\alpha,\ n-x+\beta)}{B(\alpha,\beta)} \\
        &=\binom{n}{x}\frac{\Gamma(x+\alpha)\Gamma(n-x+\beta)}{\Gamma(n+\alpha+\beta)}\cdot \frac{\Gamma(\alpha+\beta)}{\Gamma(\alpha)\Gamma(\beta)},　x=0,1,\dots,n.
        \end{align*}
        It is the Beta--Binomial pmf.

        \item Explain why mixtures increase variance.
        
        Sol.

        混合分配會使變異數增加，乃至於 Law of total variance ：
        $$Var(X)=E\left[Var(X\mid \theta)\right]+Var\left(E[X\mid \theta]\right).$$
        其中第一項是「抽樣變異」（Sampling variation），描述在參數 $\theta$ 固定時的隨機變動；
        第二項是「異質性變異」（Heterogeneity effect），描述在參數 $\theta$ 隨機時而引入的額外變動。

        若不混合，此時 $\theta$ 為常數，則 $E[X\mid \theta]$ 為常數，
        因此
        $$Var\left(E[X\mid \theta]\right)=0\Rightarrow Var(X)=E\left[Var(X\mid \theta)\right].$$
        若混合，此時 $\theta$ 為隨機變數且會影響 $E[X\mid\theta]$，則
        $$Var\left(E[X\mid \theta]\right)>0\Rightarrow Var(X)=E\left[Var(X\mid \theta)\right]+Var\left(E[X\mid \theta]\right)>E\!\left[Var(X\mid \theta)\right].$$
        因此混合會增加異質性變異，使得總變異數膨脹（Overdispersion）。

        以 Problem 4 的 Poisson–Gamma mixture 為例

        已知 $$E[X \mid \Lambda]=Var(X \mid \Lambda)=\Lambda,　E[\Lambda]=\frac{\alpha}{\beta},　Var(\Lambda)=\frac{\alpha}{\beta^{2}}$$
        可得
        \begin{align*}
            Var(X)
            &=E\left[Var(X\mid \Lambda)\right]+Var\left(E[X\mid \Lambda]\right) \\
            &=E[\Lambda]+Var(\Lambda) \\
            &=\frac{\alpha}{\beta}+\frac{\alpha}{\beta^{2}}
        \end{align*}
        我們知道
        $$E[X]=E\big(E[X \mid \Lambda]\big)=E[\Lambda]=\frac{\alpha}{\beta}.$$
        因此
        $$Var(X)=E[X]+\frac{\alpha}{\beta^{2}}>E[X].$$
        總結來說，對於一個 Poisson 分配，其變異數等於期望值，但經過 Poisso-Gamma 的混和分配後，造成變異數大於期望值，亦即變異數膨脹。
    \end{enumerate}

    \textbf{SPC Link}: Overdispersion violates p-chart assumptions.

    Sol.

    在SPC中，c-chart假設每個檢驗單位的缺陷數服從 Poisson 分配，其平均數與變異數皆為 $\lambda$ ，若缺陷發生率在不同時間或批次之間產生變動，則會形成 Poisson-Gamma 的混和分配，
    使得分配變為 Negative-Bionimal ，且變異數會大於平均數，即 Overdispersion ，因此違反 c-chart 的 Poisson 假設，使得控制界線仍在使用相對較小的變異數，導致控制界線過窄。
\end{document}